\documentclass[11pt]{article}

\usepackage[margin=1in]{geometry}
\usepackage{authblk}
\usepackage{hyperref}
\usepackage{parskip}

\title{A Triangulated, Date-Corrected Dataset of U.S. National Flood
Insurance Program Claims}

\author[1,*]{Yuya Kawakami}
% Additional author(s) -- TBD. Confirm whether Tom Corringham (basis for
% the inflation-adjustment approach in adjust_inflation.py) should be
% listed as a co-author here or credited in Acknowledgements instead.

\affil[1]{[Department], University of California, Davis, Davis, CA, USA}
\affil[*]{Corresponding author: yuya737@gmail.com}

\date{}

\begin{document}

\maketitle

\section*{Abstract}
% Draft -- 150-200 words, no citations, no subheadings, per Scientific
% Data style. Revise once Data Overview numbers below are filled in.

The National Flood Insurance Program (NFIP) Redacted Claims dataset is a
primary public record of flood losses in the United States, but two
properties of the raw extract limit its direct use in spatial and
temporal analysis. First, reported locations are deliberately coarsened
for policyholder privacy, leaving each claim with three independent,
imprecise location signals (a census block group, a ZIP code, and a
latitude/longitude rounded to one decimal place) rather than a single
precise point. Second, a subset of claims attributed to rain-driven
(pluvial) flooding carry a reported date of loss that does not coincide
with any local precipitation event, suggesting reporting or processing
error. We present a cleaning pipeline and derived dataset that (1)
adjusts all monetary fields to constant-year dollars, (2) triangulates a
bounded spatial-uncertainty polygon per claim from the three location
signals using an empirically validated, decade-based Census boundary
assignment, and (3) corrects likely-erroneous dates for pluvial claims
against NOAA's Analysis of Record for Calibration (AORC) precipitation
reanalysis, at each claim's own local time. The pipeline is fully
scripted and reproducible from public sources, and its methodological
choices are validated empirically against the claims data itself rather
than assumed.

\section*{Background \& Summary}
% Draft. Expand with citations to prior NFIP-claims literature and
% flood-risk studies that motivate this work -- none included yet.

The NFIP, administered by FEMA, is the principal source of insured flood
loss data in the United States and is widely used in flood-risk,
climate-adaptation, and disaster-economics research. FEMA releases a
redacted extract of individual claims records (OpenFEMA's ``FIMA NFIP
Redacted Claims v2'') that is public, free of authentication, and
regularly updated, but it is a \emph{redacted} extract: several
properties of the raw records make it unsuitable for direct use in
location- or date-sensitive analysis without additional processing.

Location is redacted to protect policyholder privacy. No street address
is provided. Instead, each record carries a census block group FIPS
code, a ZIP code, and a latitude/longitude pair rounded to one decimal
place --- three coarse, independently imprecise proxies for the claim's
true location, none of which alone bounds that location tightly. A
na\"ive analysis using any single one of these fields either over- or
under-states the spatial precision actually available.

Reported dates of loss are also imperfect for a subset of claims. Claims
attributed to pluvial (rain-driven) flooding are recorded in NFIP using
\texttt{causeOfDamage == "4"} (``Other''), a broad catch-all rather than
a clean pluvial label, and the reported \texttt{dateOfLoss} for such
claims does not always coincide with a day of meaningful precipitation
at the claim's location --- consistent with reporting delay, adjuster
processing lag, or transcription error rather than the actual date of
the flood event.

This paper documents an open, reproducible pipeline that addresses both
problems using only public data: FEMA's own claims extract, the Federal
Reserve Economic Data (FRED) personal consumption expenditures price
index, U.S. Census Bureau cartographic boundary files, and NOAA's AORC
v1.1 precipitation reanalysis. Every non-trivial methodological choice
in the pipeline --- which Census boundary vintage to assign to a claim,
whether time zone and daylight-saving effects matter for date
correction, what AORC grid resolution is actually available --- was
tested empirically against the claims data and the source archives
themselves rather than assumed, and the resulting evidence is reported
alongside the choice it justifies (see Methods and Technical
Validation). The output is a claim-level dataset carrying an added
spatial-uncertainty polygon and, where applicable, a corrected date of
loss, released under an open license alongside the pipeline code that
produced it.

\section*{Methods}

This pipeline processes FEMA's public claims extract in three stages,
run in order: inflation adjustment, spatial triangulation, and pluvial
date correction. Each stage is implemented as an independent,
re-runnable script; the full sequence is in \texttt{src/run\_pipeline.sh}
of the code repository (see Code Availability).

\subsection*{Input data}

\begin{itemize}
  \item \textbf{FEMA OpenFEMA} --- FIMA NFIP Redacted Claims v2
    (\texttt{FimaNfipClaims.parquet}). Public bulk download, no
    authentication. US government work, public domain.
  \item \textbf{FRED} (Federal Reserve Bank of St. Louis) --- Personal
    consumption expenditures price index, series
    \texttt{DPCERD3Q086SBEA}. Public CSV endpoint, no authentication.
  \item \textbf{U.S. Census Bureau} --- Cartographic boundary
    shapefiles: block groups and ZIP Code Tabulation Areas (ZCTAs),
    assigned per claim decade rather than one file per claim year (see
    below). Public bulk download, public domain.
  \item \textbf{NOAA} --- Analysis of Record for Calibration (AORC)
    v1.1, 30 arc-second ($\sim$800 m) grid, hourly. Public,
    \texttt{s3://noaa-nws-aorc-v1-1-1km}, no-sign-request.
\end{itemize}

None of these sources are redistributed in raw form in the code
repository; each pipeline step downloads or expects a locally staged
copy of its input, documented at the point of use.

\textbf{Scope.} This dataset covers the 48 contiguous U.S. states and
the District of Columbia (CONUS). NFIP claims also exist for Alaska,
Hawaii, Puerto Rico, and other territories, which are excluded by
design --- Census cartographic boundary conventions and AORC's grid
extent both treat CONUS as a distinct product from these territories ---
rather than dropped as an incidental side effect of processing.

\subsection*{Stage 1: Inflation adjustment}

Every dollar-denominated claim field (paid amounts, coverage limits,
building/contents value, replacement cost) is rescaled to a fixed target
year's dollars using the FRED PCE price index, matched to each claim's
loss quarter. The default target year is 2021 and is configurable.

\subsection*{Stage 2: Spatial triangulation}

Each claim ships with three independent, imprecise location signals: a
census block group FIPS code, a ZIP code, and a rounded
latitude/longitude. Triangulation intersects whichever of these are
available into a single polygon per claim, placing a strict upper bound
on the claim's true location that is typically substantially smaller
than any one source's region alone.

\textbf{Boundary vintage assignment.} Census block-group and ZCTA
boundaries are redrawn each decennial cycle, so the default assumption
would be to match each claim to the boundary vintage in effect at its
loss year. Testing against the claims data does not support that
assumption: matching \texttt{censusBlockGroupFips} against the official
2010 decennial block-group release yields a 92--100\% match rate for
claims in every decade back to the 1970s, including years before block
groups existed as a nationwide census geography. This indicates FEMA's
geocoding assigns a recent-vintage boundary to each record at processing
time rather than using the historical loss-year boundary, and there is
no sharp cutover to the newer (2020) vintage: 77.7\% of block-group
GEOIDs are identical between the 2010 and 2020 releases, and the match
rate against 2020 increases gradually with loss year rather than
switching at a specific point.

A per-claim membership test against every available shapefile release
would recover more matches than any decade-based rule (measured at
approximately 88,000 additional CONUS claims, 3.4\%, that match only the
release a decade-based table does not assign to their era), since the
transition between releases is gradual. This pipeline instead uses an
explicit per-decade shapefile mapping: each claim's decade
(\texttt{yearOfLoss // 10 * 10}) is looked up in a configuration table
that maps directly to one shapefile, with no fallback to another release
on a miss. This is a deliberate precision/auditability tradeoff,
favoring a rule that can be stated and audited in a few lines over a
data-driven per-claim selection with no documented counterpart in FEMA's
own process. Most decades are assigned the official 2010 Census release,
with only the 2020s decade assigned to the 2020 release. Together these
leave 1,017 CONUS claims (0.04\%) matching neither release, a residual
accepted without further investigation.

Block-group and ZIP lookups use the same decade key but independent
configuration tables, so they need not resolve to the same underlying
release for a given claim. ZCTAs exist only from the 2000 census onward
(verified against Census's \texttt{PREVGENZ} archive, which contains no
pre-2000 ZIP-tabulation product), so pre-2000 claims skip the ZIP source
entirely and fall back to block group and lat/lon only. ZCTA-release
choice was found to have little effect on match rate (a strict
per-decade rule loses only about 1,900 claims relative to testing all
three available ZCTA releases), unlike the block-group case, but each
decade from 2000 onward is still assigned its own contemporaneous ZCTA
release for methodological consistency.

Once a decade's shapefile is selected, the match is validated against
FEMA's rounded lat/lon rather than accepted on GEOID string match alone:
a code that matches the decade's shapefile GEOID but whose polygon does
not overlap the $0.1 \times 0.1$ degree box implied by the rounded
coordinates is treated as spatially inconsistent and excluded from the
claim's intersection. A claim's location sources can also each
individually validate against the lat/lon box while failing to overlap
one another; this is flagged explicitly rather than allowing a zero-area
polygon to be published silently.

\subsection*{Stage 3: Pluvial date correction}

For each claim coded \texttt{causeOfDamage == "4"}, the reported
\texttt{dateOfLoss} is checked against whether it coincides with a day
of meaningful precipitation (maximum hourly precipitation exceeding a
configurable threshold, default 5.0 mm) at the claim's triangulated
location, using AORC v1.1. If it does not, a window of nearby days
(default $\pm$7) is searched for the wettest day, and if that day
clears the threshold it is treated as the corrected date. The original
\texttt{dateOfLoss} is never overwritten; the corrected value and an
explicit status code are added as new columns, letting downstream users
decide whether to apply the correction.

Day boundaries are evaluated in each claim's own local time zone ---
resolved geographically per claim location --- rather than a single
national reference clock, and daylight-saving transitions are handled
correctly (a local day spans 23 or 25 hours, not always 24, across a
spring-forward or fall-back date) rather than assuming a fixed UTC
offset. Precipitation at each claim's exact location is computed by
bilinear interpolation of the four surrounding AORC grid pixels, rather
than nearest-pixel lookup, before the daily maximum is taken.

AORC coverage for CONUS begins 1979-01-01. Pluvial claims before this
date, and claims located outside the AORC CONUS grid extent, are not
evaluated and retain their reported date. Rather than downloading AORC's
full hourly CONUS archive, the pipeline computes the substantially
smaller set of distinct (pixel, day) pairs actually referenced by any
pluvial claim's search window, and fetches only those values from the
public archive.

\section*{Data Records}
% Filenames/paths below reflect the current pipeline output location;
% update if final published data records are deposited under different
% names -- see Data Availability.

The dataset is released as a single columnar file
(\texttt{claims\_pluvial\_corrected.parquet}, Apache Parquet format) at
claim-level granularity. All original columns from the FEMA OpenFEMA
\texttt{FimaNfipClaims} extract are preserved unchanged (see FEMA's own
data dictionary for those). The columns added by this pipeline are:

\textbf{Added by inflation adjustment.} \texttt{\{field\}Real2021} for
each of 15 dollar-denominated fields (paid amounts, coverage limits,
building/contents value, replacement cost), rescaled to 2021 dollars;
\texttt{amountPaid} / \texttt{amountPaidReal2021} (sum of the three
\texttt{amountPaidOn*Claim} fields); \texttt{damageAmount} /
\texttt{damageAmountReal2021} (\texttt{buildingDamageAmount +
contentsDamageAmount}).

\textbf{Added by triangulation.} \texttt{geometry} (polygon, EPSG:5070
Albers Equal Area), the intersection of whichever of \{block-group
boundary, ZCTA boundary, lat/lon box\} were available and spatially
validated for the claim; \texttt{n\_geometry\_sources} (0--3) and
\texttt{geometry\_sources} (e.g. \texttt{"block\_group+zip+latlon"})
recording which sources contributed; \texttt{geometry\_is\_empty},
flagging claims whose sources each individually validated but do not
overlap one another; \texttt{decade\_used}, the claim's decade bucket
and lookup key into the boundary-vintage configuration (see Methods);
\texttt{block\_group\_match\_status} and \texttt{zip\_match\_status},
each one of \texttt{not\_found}, \texttt{validated},
\texttt{unvalidated\_no\_latlon}, \texttt{spatially\_inconsistent}, or
\texttt{no\_shapefile\_for\_decade}.

\textbf{Added by pluvial date correction.}
\texttt{correctedDateOfLoss}, the reported date if it already coincided
with meaningful precipitation, the corrected date if a correction was
made, or the original date otherwise (always populated);
\texttt{pluvialCorrectionStatus}, one of \texttt{not\_pluvial},
\texttt{before\_aorc\_coverage}, \texttt{outside\_aorc\_grid\_extent},
\texttt{accepted\_as\_reported}, \texttt{corrected},
\texttt{no\_qualifying\_precip\_in\_window}, or
\texttt{no\_aorc\_data\_in\_window};
\texttt{pluvialCorrectionMaxPrecipMm}, the maximum daily precipitation
(mm) found at the claim's location on the date
\texttt{correctedDateOfLoss} refers to.

\texttt{geometry} uses EPSG:5070 (meters); the original FEMA
\texttt{latitude} / \texttt{longitude} fields remain in WGS84
(EPSG:4326) degrees, unchanged.

\section*{Data Overview}
% TBD -- populate once a full pipeline run against the current codebase
% (per-decade shapefile configuration and bilinear AORC interpolation)
% has been completed and verified. Include: total claim count in the
% final dataset; count and share of claims by n_geometry_sources; count
% and share of pluvial claims by pluvialCorrectionStatus, in particular
% the share actually corrected; temporal coverage (min/max yearOfLoss)
% and its distribution; and a map or summary of claim counts by state.

[TBD]

\section*{Technical Validation}

Each non-default methodological choice in this pipeline was validated
against either the claims data itself or the source archives directly,
rather than assumed from documentation alone.

\textbf{Boundary vintage assignment (Stage 2).} The decade-based
assignment rule was chosen by directly measuring match rates of
\texttt{censusBlockGroupFips} against the 2010 and 2020 official
block-group releases, by decade of \texttt{yearOfLoss}, across the full
claims extract. This measurement is what established both that a
2010-vintage release matches 92--100\% of claims in every decade back to
the 1970s (evidence that FEMA's geocoding is not
historical-vintage-aware) and that the 2010-versus-2020 transition is
gradual (77.7\% of GEOIDs identical between releases) rather than sharp
--- the empirical basis for the per-decade rule described in Methods,
including its measured 88,000-claim (3.4\%) cost relative to a per-claim
membership test. The pre-2000 ZCTA gap was verified directly against
Census's \texttt{PREVGENZ} archive (no pre-2000 ZIP-tabulation product
exists) rather than assumed from general knowledge of Census geography
changes; ZCTA-vintage sensitivity was measured directly (approximately
1,900 claims) rather than assumed negligible.

\textbf{AORC grid resolution (Stage 3).} AORC's public S3 bucket is
named \texttt{noaa-nws-aorc-v1-1-1km}, suggesting a 1 km grid. The
actual grid spacing was verified directly against the archive's
coordinate metadata: a 30 arc-second (1/120 degree) step, corresponding
to approximately 926 m north-south and 610--840 m east-west across
CONUS latitudes (varying with latitude because a degree of longitude
spans less ground distance at higher latitudes), not 1 km. This grid
resolution is stated correctly in Methods rather than inherited from the
bucket name.

\textbf{Time zone and daylight-saving handling (Stage 3).} Local-day
boundaries for AORC lookups are resolved through a two-step process: a
geographic time-zone lookup at each claim's location
(\texttt{timezonefinder}) followed by daylight-saving-aware UTC
conversion for the specific calendar date in question (Python's
\texttt{zoneinfo}), rather than a fixed UTC offset per location. This
was confirmed by direct code inspection to correctly produce 23- or
25-hour local days across spring-forward and fall-back transitions, and
to correctly attribute a December 31 local day's final hours to UTC's
following calendar year rather than dropping them.

\section*{Usage Notes}

\texttt{geometry} is a strict upper bound on a claim's true location,
not the location itself --- it is the intersection of multiple
imprecise sources, and its area varies claim to claim depending on how
many sources were available and how tightly they constrain each other
(\texttt{n\_geometry\_sources}, \texttt{geometry\_sources}). Users
requiring a fixed minimum precision should filter on these columns.

\texttt{correctedDateOfLoss} is a suggested correction, not a
replacement for \texttt{dateOfLoss}: both are retained, and
\texttt{pluvialCorrectionStatus} records why a claim was or was not
corrected, so that users can choose their own threshold for accepting
the correction or fall back to the original reported date.

This dataset is CONUS-only; NFIP claims for Alaska, Hawaii, Puerto Rico,
and other territories are not included (see Methods, Scope).

\texttt{causeOfDamage == "4"} is used as the best available proxy for
pluvial claims but is a coarse label (see Methods, Stage 3); results
that depend on precisely isolating pluvial-caused losses should treat
\texttt{pluvialCorrectionStatus == "not\_pluvial"} accordingly rather
than as a guarantee of non-pluvial cause.

An optional structure-refinement stage against FEMA's National Structure
Inventory (\texttt{refine\_with\_nsi.py}) exists in the code repository
but is not part of the pipeline run that produces the released dataset.

\section*{Data Availability}
% TBD -- fill in the permanent archive/DOI once the derived dataset is
% deposited, per Scientific Data's requirement that data be hosted in a
% recognized public repository, not solely on GitHub.

The derived dataset described in this paper is released under a CC-BY
4.0 license. [Repository name/DOI TBD.] The three upstream sources it is
built from are also public: FEMA's OpenFEMA FIMA NFIP Redacted Claims v2
(\url{https://www.fema.gov/openfema-data-page/fima-nfip-redacted-claims-v2}),
U.S. Census Bureau cartographic boundary files
(\url{https://www.census.gov/geographies/mapping-files/time-series/geo/carto-boundary-file.html}),
and NOAA's AORC v1.1 (public S3, \texttt{s3://noaa-nws-aorc-v1-1-1km},
no-sign-request). None require redistribution restrictions beyond those
already stated.

\section*{Code Availability}
% TBD -- fill in the public code repository URL once available.

The pipeline code that produces this dataset from the public sources
listed above is released under the MIT License at [repository URL TBD].
The repository includes a full methods write-up
(\texttt{docs/methods.md}) and data dictionary
(\texttt{docs/data\_dictionary.md}) documenting every column this
pipeline adds, and a single entry-point script
(\texttt{src/run\_pipeline.sh}) that reproduces the released dataset end
to end from public sources with no manual steps beyond supplying local
paths for downloaded Census shapefiles.

\section*{References}
% TBD -- add citations for: prior NFIP-claims research this work builds
% on or is intended to support; FEMA OpenFEMA data documentation; FRED
% PCE series documentation; Census cartographic boundary file
% documentation; NOAA AORC v1.1 technical documentation; timezonefinder
% if a citable reference exists.

\section*{Author Contributions}
% TBD -- confirm final author list before completing this section using
% CRediT roles.

Y.K.: conceptualization, methodology, software, validation, writing ---
original draft. [Additional authors: TBD.]

\section*{Competing Interests}

The authors declare no competing interests. [Confirm before submission.]

\section*{Acknowledgements}
% TBD. If Tom Corringham is not listed as a co-author, acknowledge here
% the notebook this pipeline's inflation-adjustment approach is adapted
% from.

\section*{Funding}
% TBD -- list any grant/award numbers and funding bodies supporting this
% work, or state that none was received.

\end{document}
